% ============================================================
%  CHAPITRE 4 : Du Modèle au Service — Mise en Production
% ============================================================
\chapter{Du Modèle au Service~: Mise en Production}
\label{chap:production}

\begin{infobox}[Résumé du chapitre]
Ce chapitre décrit le déploiement du modèle XGBoost en production~: API FastAPI,
conteneurisation Docker, pipeline CI/CD GitHub Actions, monitoring de la dérive,
rollback automatique, orchestration Prefect, analyse de fairness et bilan des
tests réalisés.
\end{infobox}

% ============================================================
\section{API de prédiction avec FastAPI}
% ============================================================

\subsection{Design de l'API}

L'API de prédiction expose trois endpoints principaux, chacun répondant à un besoin
précis du cycle de vie opérationnel. Le endpoint \texttt{/health} retourne l'état
complet du service~: version du modèle champion actif, statut du challenger en staging
et disponibilité des dépendances (MLflow, feature store). Le endpoint \texttt{/predict}
est le cœur fonctionnel~: il reçoit les 29~features d'un patient, interroge le feature
store pour appliquer les mêmes transformations de preprocessing qu'à l'entraînement,
soumet le vecteur au modèle et retourne un score de risque numérique accompagné d'un
niveau qualitatif. Le endpoint \texttt{/model/info} expose les métadonnées du modèle
en production~: version MLflow, AUC d'entraînement, liste des features attendues.
Ces métadonnées permettent aux consommateurs de l'API de détecter automatiquement
les changements de version.

\begin{table}[htbp]
\centering
\caption{Endpoints de l'API FastAPI}
\label{tab:api}
\begin{tabular}{@{}llll@{}}
\toprule
\textbf{Méthode} & \textbf{Endpoint} & \textbf{Réponse} & \textbf{SLA} \\
\midrule
GET  & \texttt{/health}      & Statut + version modèle      & $<$10ms \\
POST & \texttt{/predict}     & Score risque + niveau + seuil & $<$100ms \\
GET  & \texttt{/model/info}  & Métadonnées + AUC + features  & $<$10ms \\
\bottomrule
\end{tabular}
\end{table}

\subsection{Les trois niveaux de risque}

La réponse de \texttt{/predict} classe le score de risque en trois niveaux
cliniquement interprétables~: \textbf{FAIBLE} pour un score inférieur à 0.40,
indiquant qu'une surveillance standard suffit~; \textbf{MODÉRÉ} entre 0.40 et
0.70, suggérant un suivi renforcé par le médecin traitant~; et \textbf{ÉLEVÉ}
au-delà de 0.70, nécessitant une intervention préventive coordonnée. Le seuil
de décision final $\tau^* = 0.87$ est appliqué lors de l'interprétation binaire,
mais le score numérique continu est toujours fourni pour permettre aux cliniciens
d'adapter leur jugement.

L'architecture de l'API est conçue pour que le modèle soit chargé une seule fois
au démarrage du conteneur~: cette décision a permis de réduire la latence de 320ms
(avec chargement à chaque requête) à moins de 100ms. L'implémentation complète du
endpoint \texttt{/predict} et du schéma Pydantic des 29~features est disponible
en Annexe~A.

\begin{figure}[htbp]
\centering
\IfFileExists{figures/swagger.png}{%
  \includegraphics[width=0.9\linewidth]{figures/swagger.png}%
}{%
  \begin{tikzpicture}
    \draw[fill=lightblue, draw=primary, rounded corners=4pt]
      (0,0) rectangle (12,3.5);
    \node[font=\large\bfseries, text=primary] at (6,1.75)
      {[Capture Swagger /docs + résultat /predict]};
    \node[font=\small, text=gray80] at (6,0.75)
      {figures/swagger.png --- à remplacer};
  \end{tikzpicture}%
}
\caption{Interface Swagger --- documentation de l'API et exemple de prédiction}
\label{fig:swagger_ch4}
\end{figure}


% ============================================================
\section{Conteneurisation avec Docker}
% ============================================================

\subsection{Architecture Docker Compose}

L'ensemble du système est orchestré par Docker Compose, qui gère quatre services
interdépendants avec des règles de démarrage strictes. MLflow démarre en premier,
car l'API FastAPI en dépend pour charger le modèle de production au démarrage.
L'API et Airflow ne démarrent qu'après que MLflow soit déclaré \textit{healthy}
(sonde \texttt{healthcheck} sur \texttt{/health}). Evidently et Grafana complètent
la stack de monitoring, avec des volumes montés pour persister les dashboards
et les rapports de dérive entre les redémarrages.

Cette architecture de démarrage ordonné évite une catégorie entière de bugs de
production~: si l'API démarrait avant MLflow, le chargement du modèle échouerait
silencieusement et toutes les requêtes retourneraient une erreur 500 jusqu'au
prochain redémarrage manuel. La commande de démarrage complète (\texttt{docker
compose up -d}) ainsi que les commandes de diagnostic sont disponibles en Annexe~A.

\begin{figure}[htbp]
\centering
\IfFileExists{figures/docker.png}{%
  \includegraphics[width=0.9\linewidth]{figures/docker.png}%
}{%
  \begin{tikzpicture}
    \draw[fill=lightblue, draw=primary, rounded corners=4pt]
      (0,0) rectangle (12,3.5);
    \node[font=\large\bfseries, text=primary] at (6,1.75)
      {[Capture conteneurs Docker actifs]};
    \node[font=\small, text=gray80] at (6,0.75)
      {figures/docker.png --- à remplacer};
  \end{tikzpicture}%
}
\caption{Conteneurs Docker actifs --- les quatre services démarrés avec succès}
\label{fig:docker}
\end{figure}

\subsection{Isolation et reproductibilité}

Chaque service dispose de son propre Dockerfile avec des versions épinglées
(\textit{pinned}) pour toutes les dépendances. Cette pratique garantit que
l'environnement de production est identique à l'environnement de développement,
éliminant les problèmes classiques du type \og ça marche sur ma machine\fg{}. Les
images Docker construites sont taguées avec le hash du commit Git correspondant,
créant une traçabilité complète entre le code, l'image et les prédictions produites.

% ============================================================
\section{CI/CD avec GitHub Actions}
% ============================================================

\subsection{Pipeline d'intégration continue}

Le pipeline CI/CD déclenché à chaque \textit{push} sur la branche principale
exécute quatre étapes séquentielles. La première valide les données via Great
Expectations~: si les données en entrée ne respectent pas les contrats définis,
le pipeline s'arrête et aucun déploiement n'a lieu. La deuxième exécute la
suite de tests unitaires et d'intégration avec \texttt{pytest}~: tests de
l'API (codes HTTP, format des réponses), tests des transformations de features
et tests de chargement du modèle. La troisième construit l'image Docker et la
publie dans le registre de conteneurs GitHub. La quatrième redéploie les
services en production avec \texttt{docker compose up -d --force-recreate}.

Ce pipeline garantit que toute modification du code, même minime, passe par
une validation complète avant d'atteindre la production. La configuration
complète du fichier \texttt{.github/workflows/ci.yml} est disponible en Annexe~A.

\begin{alertbox}[Gestion des secrets CI/CD]
Les credentials MLflow, les clés d'accès S3 et les tokens GitHub sont stockés
dans les \textit{GitHub Secrets} et injectés comme variables d'environnement
lors de l'exécution du pipeline. Aucune donnée sensible n'est présente en clair
dans le dépôt.
\end{alertbox}

% ============================================================
\section{Monitoring~: détecter la dérive avant qu'elle ne dégrade}
% ============================================================

\subsection{Seuils d'alerte}

Le monitoring repose sur cinq métriques surveillées en continu, chacune avec
un seuil d'alerte calibré sur les caractéristiques du projet. Le PSI
(\textit{Population Stability Index}) mesure la dérive des features entre la
distribution de référence (données d'entraînement 2008) et les données de
production courantes. Un PSI supérieur à 0.20 déclenche automatiquement un
cycle de réentraînement. L'AUC en production est estimée périodiquement par
Evidently AI à partir d'un sous-ensemble de données labellisées rétrospectivement~;
si elle tombe sous 0.75, le rollback est déclenché immédiatement.

\begin{table}[htbp]
\centering
\caption{Seuils de monitoring et actions déclenchées}
\label{tab:monitoring}
\begin{tabular}{@{}lll@{}}
\toprule
\textbf{Métrique} & \textbf{Seuil} & \textbf{Action} \\
\midrule
PSI données          & $>~0.20$                  & Ré-entraînement automatique \\
AUC production       & $<~0.75$                  & Rollback automatique \\
Dégradation AUC      & $>~5\,\%$ sur 2 semaines  & Alerte équipe \\
Recall               & $<~0.70$                  & Ré-entraînement \\
Latence P99          & $>~150$~ms                & Alerte infrastructure \\
\bottomrule
\end{tabular}
\end{table}

\subsection{Simulation de dérive 2010 → 2026}

Pour valider le système de monitoring sans attendre qu'une dérive réelle se
produise, une dérive artificielle a été simulée sur les données de test 2010,
mimant les changements démographiques et économiques attendus d'ici 2026~: hausse
des coûts médicaux de 30\,\%, vieillissement de la population de 3~ans en moyenne,
augmentation du nombre de comorbidités liée à une meilleure survie des patients
chroniques, et progression du taux d'hospitalisation sous l'effet du vieillissement.

\begin{table}[htbp]
\centering
\caption{Simulation de dérive temporelle 2010 → 2026}
\label{tab:drift}
\begin{tabular}{@{}llll@{}}
\toprule
\textbf{Variable} & \textbf{Perturbation} & \textbf{PSI} & \textbf{Alerte} \\
\midrule
\texttt{COUT\_TOTAL}    & $+30\,\%$          & 0.28 & \textcolor{danger}{\textbf{Oui}} \\
\texttt{AGE}            & $+3$~ans           & 0.19 & Non \\
\texttt{NB\_COMORBIDITES} & $+0.5$           & 0.31 & \textcolor{danger}{\textbf{Oui}} \\
Taux hosp.              & 7.5\,\% → 12\,\%  & 0.44 & \textcolor{danger}{\textbf{Oui}} \\
\bottomrule
\end{tabular}
\end{table}

Trois des quatre variables testées dépassent le seuil PSI~$=~0.20$, confirmant
que le monitoring détecte correctement des changements de population représentatifs
d'un glissement de 16~ans. Seul le vieillissement de 3~ans (PSI~=~0.19) passe
juste en dessous du seuil~: ce résultat est cohérent avec la nature progressive
du vieillissement, qui ne modifie pas brutalement la distribution de l'âge.

\subsection{Boucle de rollback automatique}

\begin{figure}[htbp]
\centering
\begin{tikzpicture}[
  font=\small, >=Stealth,
  node distance=1.5cm,
  box/.style={rectangle, rounded corners=4pt, minimum width=3.5cm,
              minimum height=0.9cm, text centered, draw, thick},
  decision/.style={diamond, minimum width=2.8cm, minimum height=0.9cm,
                   text centered, draw=accent, thick, fill=lightyellow},
  action/.style={box, fill=lightgreen, draw=success},
  alert/.style={box, fill=danger!10, draw=danger},
  normal/.style={box, fill=lightblue, draw=primary},
]
\node[normal]   (collect) at (0,0)    {Collecte métriques\\(Evidently)};
\node[decision] (psi)     at (0,-2.2) {PSI $>$ 0.20\\ou AUC $<$ 0.75~?};
\node[action]   (ok)      at (-4,-4)  {Système OK\\Continuer monitoring};
\node[alert]    (alert)   at (4,-4)   {Alerte Grafana\\+ Log MLflow};
\node[action]   (rollback)at (4,-6.2) {Rollback vers\\version précédente};
\node[normal]   (retrain) at (4,-8.4) {Ré-entraînement\\planifié (Prefect)};
\draw[->, thick] (collect) -- (psi);
\draw[->, thick] (psi) -- node[left, font=\footnotesize] {Non} (ok);
\draw[->, thick] (psi) -- node[right, font=\footnotesize] {Oui} (alert);
\draw[->, thick] (alert) -- (rollback);
\draw[->, thick] (rollback) -- (retrain);
\draw[->, dashed, draw=primary] (retrain) -- ++(3,0) -- ++(0,8.4) -- (collect);
\end{tikzpicture}
\caption{Boucle de rollback automatique~--- détection, alerte et récupération}
\label{fig:rollback}
\end{figure}

La boucle de la figure~\ref{fig:rollback} s'exécute sans intervention humaine.
Lorsque Evidently détecte une dérive critique, Grafana émet une alerte et
l'API bascule automatiquement vers la version précédente du modèle, enregistrée
dans le registre MLflow. Ce mécanisme garantit une continuité de service~:
la dégradation est limitée à la fenêtre temporelle entre deux exécutions
quotidiennes du pipeline de monitoring.


\begin{figure}[htbp]
\centering
\IfFileExists{figures/grafana.png}{%
  \includegraphics[width=0.9\linewidth]{figures/grafana.png}%
}{%
  \begin{tikzpicture}
    \draw[fill=lightblue, draw=primary, rounded corners=4pt]
      (0,0) rectangle (12,3.5);
    \node[font=\large\bfseries, text=primary] at (6,1.75)
      {[Dashboard Grafana/Evidently]};
    \node[font=\small, text=gray80] at (6,0.75)
      {figures/grafana.png --- à remplacer};
  \end{tikzpicture}%
}
\caption{Dashboard Grafana~--- monitoring temps réel avec alertes Evidently AI}
\label{fig:grafana}
\end{figure}

\begin{figure}[htbp]
\centering
\IfFileExists{figures/evidently.png}{%
  \includegraphics[width=0.9\linewidth]{figures/evidently.png}%
}{%
  \begin{tikzpicture}
    \draw[fill=lightblue, draw=primary, rounded corners=4pt]
      (0,0) rectangle (12,3.5);
    \node[font=\large\bfseries, text=primary] at (6,1.75)
      {[Rapport Evidently drift]};
    \node[font=\small, text=gray80] at (6,0.75)
      {figures/evidently.png --- à remplacer};
  \end{tikzpicture}%
}
\caption{Rapport Evidently AI~--- détail des colonnes en dérive avec PSI par feature}
\label{fig:evidently}
\end{figure}

% ============================================================
\section{Orchestration avec Prefect}
% ============================================================

Prefect~\cite{prefect2023} orchestre les pipelines de données et de réentraînement
selon un planning déclaratif. Le flow principal de réentraînement s'articule en
quatre tâches séquentielles~: validation des nouvelles données par Great Expectations,
réentraînement sur les données historiques complètes, évaluation du nouveau modèle
et enregistrement conditionnel dans MLflow si l'AUC dépasse le seuil de 0.85.
Chaque tâche est configurée avec trois tentatives automatiques en cas d'échec,
avec un délai de 60~secondes entre chaque tentative.

L'atout principal de Prefect par rapport à Airflow pour ce pipeline est la gestion
native des artefacts~: chaque exécution produit un rapport traçable accessible via
l'interface Prefect Cloud, avec les logs de chaque tâche et les résultats des
validations. L'implémentation complète du flow Prefect est disponible en Annexe~A.

\begin{figure}[htbp]
\centering
\IfFileExists{figures/prefect.png}{%
  \includegraphics[width=0.9\linewidth]{figures/prefect.png}%
}{%
  \begin{tikzpicture}
    \draw[fill=lightblue, draw=primary, rounded corners=4pt]
      (0,0) rectangle (12,3.5);
    \node[font=\large\bfseries, text=primary] at (6,1.75)
      {[Interface Prefect]};
    \node[font=\small, text=gray80] at (6,0.75)
      {figures/prefect.png --- à remplacer};
  \end{tikzpicture}%
}
\caption{Interface Prefect~--- historique des runs et statut des tâches}
\label{fig:prefect}
\end{figure}

% ============================================================
\section{Analyse de fairness~: équité algorithmique}
% ============================================================

\subsection{Résultats par groupe racial}

L'analyse de fairness évalue si le modèle performe équitablement pour tous les
groupes de patients. Un système qui obtient une excellente AUC globale mais qui
prédit mal pour un groupe minoritaire n'est pas acceptable en contexte médical~:
il introduirait un biais systématique qui pénaliserait précisément les populations
déjà les plus vulnérables sur le plan socio-économique~\cite{obermeyer2019dissecting}.

\begin{table}[htbp]
\centering
\caption{Analyse de fairness par groupe racial}
\label{tab:fairness}
\begin{tabular}{@{}lrrrr@{}}
\toprule
\textbf{Groupe racial} & \textbf{N} & \textbf{AUC} & \textbf{Recall} & \textbf{FP rate} \\
\midrule
Blanc           & 21\,840 & 0.901 & 0.894 & 0.241 \\
Noir            &  4\,210 & 0.887 & 0.871 & 0.268 \\
Autre           &  1\,520 & 0.883 & 0.863 & 0.274 \\
Asiatique       &    830  & 0.879 & 0.858 & 0.281 \\
Hispanique      &    520  & 0.876 & 0.851 & 0.289 \\
Natif Américain &    155  & 0.871 & 0.839 & 0.295 \\
\midrule
\textbf{Écart max Recall} & & & \textbf{5.5\,\%} & \\
\bottomrule
\end{tabular}
\end{table}

\begin{alertbox}[Interprétation de l'écart de fairness]
L'écart maximal de \textbf{5.5\,\%} en rappel entre le groupe Blanc et le groupe
Natif Américain est statistiquement observable mais cliniquement modéré. Il
s'explique principalement par le déséquilibre des effectifs~: avec seulement
155~patients natifs américains, l'estimation du rappel est moins stable.
Les perspectives d'amélioration incluent~:
\begin{itemize}
  \item sur-représentation lors du ré-entraînement (\textit{oversampling}
        ciblé par groupe) ;
  \item contraintes de fairness lors de l'optimisation Optuna~\cite{chen2021ethical}.
\end{itemize}
\end{alertbox}

Ce gradient de performance~--- de 0.901 pour le groupe le plus représenté à 0.871
pour le moins représenté~--- est congruent avec les résultats de la littérature sur
les algorithmes de santé~\cite{obermeyer2019dissecting}. La traçabilité des métriques
de fairness dans MLflow permet de suivre leur évolution à chaque réentraînement et
de déclencher une alerte si l'écart s'aggrave.

% ============================================================
\section{Bilan des tests réalisés}
% ============================================================

\begin{table}[htbp]
\centering
\caption{Tests réalisés et résultats}
\label{tab:tests}
\begin{tabular}{@{}lll@{}}
\toprule
\textbf{Test} & \textbf{Résultat} & \textbf{Détail} \\
\midrule
Great Expectations     & \textcolor{success}{\textbf{2/2 ✓}} & 88\,585 lignes, features validées \\
Data Cleaning          & \textcolor{success}{\textbf{✓}} & 88\,585 lignes, 0 doublon \\
Feature Engineering    & \textcolor{success}{\textbf{✓}} & 29 features générées \\
XGBoost entraînement   & \textcolor{success}{\textbf{✓}} & AUC~0.9245 sur validation \\
API \texttt{/predict}  & \textcolor{success}{\textbf{✓}} & $<$100ms testé avec k6 \\
Docker compose up      & \textcolor{success}{\textbf{✓}} & 4 services actifs \\
Drift PSI simulé       & \textcolor{success}{\textbf{✓}} & PSI~0.31 détecté correctement \\
Rollback automatique   & \textcolor{success}{\textbf{✓}} & Basculement confirmé \\
Fairness analyse       & \textcolor{success}{\textbf{✓}} & Écart 5.5\,\% documenté \\
\bottomrule
\end{tabular}
\end{table}

% ============================================================
\section{Difficultés rencontrées et solutions}
% ============================================================

Le développement de cette plateforme a rencontré cinq difficultés techniques majeures
qui méritent d'être documentées pour tout projet similaire. La première concernait le
parsing des dates~: le format flottant \texttt{20100312.0} issu des données CMS
provoquait des erreurs silencieuses qui corrompaient les calculs d'âge et de fenêtres
temporelles. La solution~--- extraction de la partie entière avant parsing~--- paraît
triviale rétrospectivement, mais a nécessité plusieurs heures de débogage.

Le deuxième défi était l'overfitting sur les données 2008, liée à l'hétérogénéité
inter-annuelle~: le modèle apprenait des patterns spécifiques à 2008 qui ne se
généralisaient pas. L'ajout de la régularisation Optuna sur les paramètres
\texttt{subsample} et \texttt{colsample\_bytree} a résolu ce problème.

\begin{table}[htbp]
\centering
\caption{Principales difficultés rencontrées et solutions adoptées}
\label{tab:difficultes}
\begin{tabular}{@{}p{4cm}p{5cm}p{4cm}@{}}
\toprule
\textbf{Difficulté} & \textbf{Cause} & \textbf{Solution} \\
\midrule
Dates non parsables    & Format flottant \texttt{.0}
  & Extraction partie entière + validation GE \\
Déséquilibre 1:12      & Taux hosp. faible
  & \texttt{scale\_pos\_weight} + seuil optimal \\
Overfitting sur 2008   & Hétérogénéité inter-annuelle
  & Split temporel + régularisation Optuna \\
Latence initiale 320ms & Chargement modèle à chaque requête
  & Chargement unique au démarrage \\
PSI faux positifs      & Winsorisation asymétrique
  & Calcul PSI post-winsorisation \\
\bottomrule
\end{tabular}
\end{table}

\section*{Synthèse du chapitre}

Ce chapitre a démontré la mise en production complète de la plateforme~: une API
FastAPI répondant en moins de 100ms, un système Docker Compose avec quatre services,
un pipeline CI/CD qui déploie automatiquement à chaque commit et un monitoring qui
détecte et corrige la dérive sans intervention humaine. L'analyse de fairness confirme
que le modèle est équitable avec un écart maximal de 5.5\,\% entre groupes raciaux,
documenté et traçable dans MLflow. La simulation de dérive 2010~→~2026 valide la
robustesse du système face à des changements démographiques réalistes.

